% ============================== 6. ИНФОРМАЦИОННИ ИЗТОЧНИЦИ ==============================
\chapter{Информационни източници}\label{ch:refs}
\begingroup
\makeatletter
\renewcommand{\chapter}{\@ifstar\@gobble\@gobble}
\makeatother
\small
\begin{thebibliography}{99}\setlength{\itemsep}{3pt}
\bibitem{antonakakis2017} Antonakakis, M. et al. Understanding the Mirai Botnet. В: \emph{26th USENIX Security Symposium}, с.~1093--1110, 2017.
\bibitem{sommer2010} Sommer, R., Paxson, V. Outside the Closed World: On Using Machine Learning for Network Intrusion Detection. В: \emph{IEEE Symposium on Security and Privacy}, с.~305--316, 2010.
\bibitem{buczak2016} Buczak, A.~L., Guven, E. A Survey of Data Mining and Machine Learning Methods for Cyber Security Intrusion Detection. \emph{IEEE Communications Surveys \& Tutorials}, 18(2):1153--1176, 2016.
\bibitem{arp2022} Arp, D. et al. Dos and Don'ts of Machine Learning in Computer Security. В: \emph{31st USENIX Security Symposium}, с.~3971--3988, 2022.
\bibitem{mirkovic2004} Mirkovic, J., Reiher, P. A Taxonomy of DDoS Attack and DDoS Defense Mechanisms. \emph{ACM SIGCOMM Computer Communication Review}, 34(2):39--53, 2004.
\bibitem{sharafaldin2018} Sharafaldin, I., Lashkari, A.~H., Ghorbani, A.~A. Toward Generating a New Intrusion Detection Dataset and Intrusion Traffic Characterization. В: \emph{ICISSP}, с.~108--116, 2018.
\bibitem{tavallaee2009} Tavallaee, M., Bagheri, E., Lu, W., Ghorbani, A.~A. A Detailed Analysis of the KDD CUP 99 Data Set. В: \emph{IEEE Symposium on Computational Intelligence for Security and Defense Applications}, 2009.
\bibitem{cse2018web} Canadian Institute for Cybersecurity. A Realistic Cyber Defense Dataset (CSE-CIC-IDS2018). University of New Brunswick. \url{https://www.unb.ca/cic/datasets/ids-2018.html} (посетен на 13.11.2026).
\bibitem{cic2017web} Canadian Institute for Cybersecurity. Intrusion Detection Evaluation Dataset (CIC-IDS2017). University of New Brunswick. \url{https://www.unb.ca/cic/datasets/ids-2017.html} (посетен на 15.11.2026).
\bibitem{neto2023} Neto, E.~C.~P., Dadkhah, S., Ferreira, R., Zohourian, A., Lu, R., Ghorbani, A.~A. CICIoT2023: A Real-Time Dataset and Benchmark for Large-Scale Attacks in IoT Environment. \emph{Sensors}, 23(13):5941, 2023. \url{https://doi.org/10.3390/s23135941}
\bibitem{cic2023web} Canadian Institute for Cybersecurity. CIC IoT Dataset 2023. University of New Brunswick. \url{https://www.unb.ca/cic/datasets/iotdataset-2023.html} (посетен на 06.11.2026).
\bibitem{lyon2009} Lyon, G.~F. \emph{Nmap Network Scanning: The Official Nmap Project Guide to Network Discovery and Security Scanning}. Insecure.Com LLC, 2009.
\bibitem{aws2018} Registry of Open Data on AWS. A Realistic Cyber Defense Dataset (CSE-CIC-IDS2018). \url{https://registry.opendata.aws/cse-cic-ids2018/} (посетен на 13.11.2026).
\bibitem{leevy2020} Leevy, J.~L., Khoshgoftaar, T.~M. A Survey and Analysis of Intrusion Detection Models Based on CSE-CIC-IDS2018 Big Data. \emph{Journal of Big Data}, 7:104, 2020.
\bibitem{engelen2021} Engelen, G., Rimmer, V., Joosen, W. Troubleshooting an Intrusion Detection Dataset: the CICIDS2017 Case Study. В: \emph{IEEE Security and Privacy Workshops}, с.~7--12, 2021.
\bibitem{liu2022} Liu, L., Engelen, G., Lynar, T., Essam, D., Joosen, W. Error Prevalence in NIDS Datasets: A Case Study on CIC-IDS-2017 and CSE-CIC-IDS-2018. В: \emph{IEEE Conference on Communications and Network Security (CNS)}, с.~254--262, 2022.
\bibitem{liu2022web} Liu, L., Engelen, G., Lynar, T., Essam, D., Joosen, W. Improved CIC-IDS-2017 and CSE-CIC-IDS-2018 Datasets. DistriNet, KU Leuven. \url{https://intrusion-detection.distrinet-research.be/CNS2022/} (посетен на 15.11.2026).
\bibitem{borovska2022} Боровска, П. Метрика за оценка на качеството на моделите и правилата за машинно обучение. Учебен материал към курса „Интелигентни системи за сигурност на екосистемите на големите данни и Интернет на нещата“, 2022.
\bibitem{guyon2003} Guyon, I., Elisseeff, A. An Introduction to Variable and Feature Selection. \emph{Journal of Machine Learning Research}, 3:1157--1182, 2003.
\bibitem{jolliffe2016} Jolliffe, I.~T., Cadima, J. Principal Component Analysis: a Review and Recent Developments. \emph{Philosophical Transactions of the Royal Society A}, 374:20150202, 2016.
\bibitem{pedregosa2011} Pedregosa, F. et al. Scikit-learn: Machine Learning in Python. \emph{Journal of Machine Learning Research}, 12:2825--2830, 2011.
\bibitem{hastie2009} Hastie, T., Tibshirani, R., Friedman, J. \emph{The Elements of Statistical Learning}. 2nd ed., Springer, 2009.
\bibitem{breiman1984} Breiman, L., Friedman, J., Olshen, R., Stone, C. \emph{Classification and Regression Trees}. Wadsworth, 1984.
\bibitem{breiman2001} Breiman, L. Random Forests. \emph{Machine Learning}, 45(1):5--32, 2001.
\bibitem{chen2016} Chen, T., Guestrin, C. XGBoost: A Scalable Tree Boosting System. В: \emph{Proceedings of the 22nd ACM SIGKDD International Conference on Knowledge Discovery and Data Mining}, с.~785--794, 2016.
\bibitem{friedman2001} Friedman, J.~H. Greedy Function Approximation: A Gradient Boosting Machine. \emph{Annals of Statistics}, 29(5):1189--1232, 2001.
\bibitem{cortes1995} Cortes, C., Vapnik, V. Support-Vector Networks. \emph{Machine Learning}, 20(3):273--297, 1995.
\bibitem{scholkopf2002} Schölkopf, B., Smola, A.~J. \emph{Learning with Kernels}. MIT Press, 2002.
\bibitem{chang2011} Chang, C.-C., Lin, C.-J. LIBSVM: A Library for Support Vector Machines. \emph{ACM Transactions on Intelligent Systems and Technology}, 2(3):27, 2011.
\bibitem{kohavi1995} Kohavi, R. A Study of Cross-Validation and Bootstrap for Accuracy Estimation and Model Selection. В: \emph{IJCAI}, с.~1137--1145, 1995.
\bibitem{fawcett2006} Fawcett, T. An Introduction to ROC Analysis. \emph{Pattern Recognition Letters}, 27(8):861--874, 2006.
\bibitem{saito2015} Saito, T., Rehmsmeier, M. The Precision-Recall Plot Is More Informative than the ROC Plot When Evaluating Binary Classifiers on Imbalanced Datasets. \emph{PLoS ONE}, 10(3):e0118432, 2015.
\bibitem{axelsson2000} Axelsson, S. The Base-Rate Fallacy and the Difficulty of Intrusion Detection. \emph{ACM Transactions on Information and System Security}, 3(3):186--205, 2000.
\bibitem{efron1993} Efron, B., Tibshirani, R.~J. \emph{An Introduction to the Bootstrap}. Chapman \& Hall, 1993.
\bibitem{mcnemar1947} McNemar, Q. Note on the Sampling Error of the Difference Between Correlated Proportions or Percentages. \emph{Psychometrika}, 12(2):153--157, 1947.
\bibitem{dietterich1998} Dietterich, T.~G. Approximate Statistical Tests for Comparing Supervised Classification Learning Algorithms. \emph{Neural Computation}, 10(7):1895--1923, 1998.
\bibitem{repo} Програмен код и резултати на изследването: хранилище \texttt{ml-ids-cicids2018} (код, конфигурация, журнали и фигури), 2026. \url{https://github.com/Vankata03/ml-ids-cicids2018} (посетен на 15.11.2026).
\bibitem{fisher2019} Fisher, A., Rudin, C., Dominici, F. All Models are Wrong, but Many are Useful: Learning a Variable's Importance by Studying an Entire Class of Prediction Models Simultaneously. \emph{Journal of Machine Learning Research}, 20(177):1--81, 2019.
\bibitem{vandermaaten2008} van der Maaten, L., Hinton, G. Visualizing Data using t-SNE. \emph{Journal of Machine Learning Research}, 9:2579--2605, 2008.
\end{thebibliography}
\endgroup
